\chapter{Adding a model}\label{ch:adding-model}

Today the engine's main model is Qwen3.8-27B, a hybrid of Gated DeltaNet and attention layers, and the small models beside it (the Qwen3.5 narrator and shadow reader) run on the same forward code. TandemLLM is meant to serve more main models than this one. This chapter lists the seams where a new model plugs in, what each one expects, and the tests a new model must pass before its numbers mean anything. Some of the engine is still specific to this family, and the last section says which parts.

\section{The seams}

\begin{figure}[htbp]
\centering
\begin{tikzpicture}
  \node[quiet, minimum width=24mm] (env) at (0,1.5) {environment};
  \node[quiet, minimum width=24mm] (cfg) at (0,0) {\texttt{config.json}};
  \node[quiet, minimum width=24mm] (ck) at (0,-1.5) {checkpoint files};
  \node[box, minimum width=28mm] (set) at (3.6,1.5) {\texttt{EngineSettings}\\[-1pt]{\scriptsize engine/settings.py}};
  \node[box, minimum width=28mm] (tc) at (3.6,0) {\texttt{TextConfig}\\[-1pt]{\scriptsize engine/config.py}};
  \node[box, minimum width=28mm] (lay) at (3.6,-1.5) {\texttt{Layout}\\[-1pt]{\scriptsize engine/loader.py}};
  \node[box, minimum width=26mm] (lin) at (7.2,-1.5) {\texttt{Linear}\\[-1pt]{\scriptsize engine/linear.py}};
  \node[hot, minimum width=28mm] (model) at (10.6,0) {target model\\[-1pt]{\scriptsize engine/model.py}};
  \node[box, minimum width=28mm] (dr) at (10.6,-2.8) {drafters declare\\[-1pt]{\scriptsize what they read}};
  \draw[flow] (env) -- (set);
  \draw[flow] (cfg) -- (tc);
  \draw[flow] (ck) -- (lay);
  \draw[flow] (lay) -- (lin);
  \draw[flow] (set.east) -| ([xshift=-4mm]model.north);
  \draw[flow] (tc) -- (model);
  \draw[flow] (lin.east) -| ([xshift=-6mm]model.south);
  \draw[flow] ([xshift=6mm]model.south) -- ([xshift=6mm]dr.north);
\end{tikzpicture}
\caption{The seams a new model plugs into.}\label{fig:seams}
\end{figure}

\paragraph{Config.} \code{engine/config.py::load_config} reads \code{config.json} into a \code{TextConfig} with every size the engine needs, from the layer count to the rotary settings, and answers which layers are recurrent. It accepts a vision-language config (the text model under \code{text_config}) and a text-only one. A new family either fits \code{TextConfig} or extends it; the rest of the engine asks \code{TextConfig}, never the raw JSON.

\paragraph{Layout.} \code{engine/loader.py::Layout} says how a checkpoint names its tensors, in one place. The engine addresses every tensor by a canonical name, the language model's own naming with the checkpoint's wrapper prefix taken off. The layout lists the prefixes to strip, the tensors the language model never reads (such as a vision tower), where the checkpoint keeps its prediction head, and the projections that may be stored as FP8 pairs or as plain matrices. \code{tests/test_layout.py} pins the served checkpoint's key set: 1,250 language-model keys read, 333 vision keys skipped.

\paragraph{Weight formats.} Every weight the engine multiplies by implements one small interface (\code{engine/linear.py}): a shape, the bytes read per use, and \code{matmul}. \code{FP8Block}, \code{NVFP4Block}, \code{FP8Head} and \code{BF16Block} implement it. The model calls \code{w.matmul(x)} and never asks which format it holds. A new format is a new class with those three members, plus a kernel that keeps the row-invariance rule (Chapter~\ref{ch:kernels}).

\paragraph{Settings.} \code{engine/settings.py} is the registry of every \code{QWEN38_*} switch, read only through \code{SETTINGS.get(name)}. \code{EngineSettings(env, prefix)} gives the same reads over another mapping or another prefix.

\paragraph{What a drafter reads.} A drafter reads parts of the target: hidden states at some layers, the embedding, the head. Each drafter declares these through \code{Drafter.requires()}: the hidden size, the tap layers, the target tensors it reads by name, and the vocabulary size. At load the engine compares them with the target and refuses a mismatch by name, so a DFlash2 drafter trained for another model fails there instead of drafting garbage. A new model needs its own block drafter; the lookup drafter works for any model once its corpus is rebuilt with the new tokenizer (\code{tools/build_corpus.py}). The lookup stores record the tokenizer's sha256 and refuse to open under another tokenizer.

\paragraph{The tokenizer and the chat template.} The server renders prompts with the model's own chat template and reads the model's own tool-call and reasoning formats. A model with another format needs its own parser behind the same OpenAI output.

\section{The tests a new model must pass}

In this order, each on the new model:

\begin{enumerate}
\item \code{tools/refcheck.py}: argmax agreement with the published reference implementation on teacher-forced positions, the reference in its own process because only one model fits on the board at a time. On Qwen3.8 the engine agrees on 99.25~\% of 400 positions and on all 242 confident ones (\measured).
\item \code{tools/quality_gate.py} for every quantised weight set (Chapter~\ref{ch:quantisation}).
\item \code{tools/verify_spec.py}\index{tools/verifyspec.py@\texttt{tools/verify\_spec.py}}, the agreement gate with its near-tie rule, with every drafter (Chapter~\ref{ch:exactness}).
\item The CPU suite, including \code{tests/test_refactor_gate.py}. That test records the exact bytes of a tiny model's run (every logit and every byte of state after a prefill, a decode, a block verify, a rollback, a tree verify and a commit) and fails if any byte moves. A change that is supposed to change arithmetic re-records the fixture and says why.
\item \code{ops/gate.sh}\index{ops/gate.sh@\texttt{ops/gate.sh}} and the client smoke set before any number is published (Chapter~\ref{ch:measurement}).
\end{enumerate}

\section{What is still specific to Qwen3.8}

\begin{itemize}
\item \code{engine/model.py} implements this family's forward: the Gated DeltaNet layer with its gated norm, attention with a gated output and partial rotary, and the dense MLP. A new family needs its own layer code behind the same verify, commit and cache interfaces; a mixture-of-experts model also needs expert routing inside the verify.
\item The recurrent-layer kernels assume Gated DeltaNet's state shape, heads of 128 by 128.
\item Tile tables and router prices are measured for this model's shapes on the GB10. A new model or a new board needs its own sweep and its own price curve.
\item Some modules read their switches into globals at import, so two engines with different settings cannot share one process yet.
\item Runtime names still carry the old project name, from the \code{QWEN38_} and \code{QSE_} prefixes to the \code{Qwen38Engine} class.
\end{itemize}

The Qwen3.5 small models already run on the engine's own forward: the narrator and the shadow reader load them through \code{engine/generative/qwen35.py}, the same layer code, because they share Qwen3.8's architecture family and its tokenizer (Chapter~\ref{ch:lines}).
